% ============================================================
%  Pi2 Final Report — PRL-style (REVTeX 4.2)
%  Target length: ~5 pages in PRL format
%
%  Compile:  latexmk -pdf main.tex   (or use the Makefile: `make`)
%  Overleaf: upload this folder as a project; it compiles as-is.
%
%  Submission: export the PDF and name it
%      Zaman_Leigh_Pi2_Report.pdf
%  (`make submission` does this automatically)
% ============================================================
\documentclass[
  aps,
  prl,
  twocolumn,
  reprint,              % journal-style layout; use `preprint` for single-column draft
  superscriptaddress,
  amsmath,
  amssymb,
  floatfix,
  longbibliography      % show article titles in the bibliography (optional)
]{revtex4-2}

\usepackage{style} % style.sty packages / macris

\begin{document}

% ============================================================
%  TITLE BLOCK
%  Guide asks for: Project Title / Scholar name / Mentor name /
%  Professor lab
% ============================================================
\title{Project Title}   % TODO: project title

\author{Leigh Zaman}
\email{leighzaman@berkeley.edu}
\affiliation{Physics Innovators Initiative (Pi$^2$), Department of Physics, University of California, Berkeley, CA 94720, USA}

% TODO: mentor name
\author{Anke Stoeltzel}
\thanks{Graduate Mentor}
\affiliation{Kolkowitz Group, Department of Physics, University of California, Berkeley, CA 94720, USA}

\date{\today}

% ============================================================
%  ABSTRACT (comes BEFORE \maketitle in REVTeX)
% ============================================================
\begin{abstract}
  % TODO: abstract (~100–150 words): context, what was done, key
  % quantitative result, significance.
\end{abstract}

\maketitle

% ============================================================
%  I. BACKGROUND AND AIMS
%  - Background and motivation of the project
%  - Major questions to address
%  - Hypothesis in addressing the question
% ============================================================
\section{Background and Aims}

% Describe the background and motivation of the project.

Identify the major questions I want to address.

Outline your hypothesis in addressing the question.

% ============================================================
%  II. EXPERIMENTAL/THEORETICAL APPROACHES
%  - Describe in detail the experimental and/or theoretical
%    approaches of the project
% ============================================================
\section{Experimental and Theoretical Approaches}

% Describe in detail the experimental and/or theoretical approaches of the project.

% Example numbered equation:
% \begin{equation}
%   I(r) \propto J_0^2(k_r r),
%   \label{eq:bessel}
% \end{equation}
% referenced as Eq.~\eqref{eq:bessel}.

% ============================================================
%  III. EXPERIMENTAL/THEORETICAL RESULTS
%  - Descriptions/figures/tables of findings
%  - New techniques and methods developed
%  - How the research addresses the posed questions
% ============================================================
\section{Results}
% TODO

% Descriptions/figures/tables of the experimental and/or theoretical findings of your project.
% Descriptions of new techniques and methods developed in your project.
% How does your research address the previously posed questions?


% Example single-column figure (put files in figures/):
% \begin{figure}[tb]
%   \centering
%   \includegraphics[width=\columnwidth]{example.pdf}
%   \caption{Caption text. \label{fig:example}}
% \end{figure}
% Referenced as Fig.~\ref{fig:example}.

% Example two-column-spanning figure:
% \begin{figure*}[tb]
%   \centering
%   \includegraphics[width=\textwidth]{example_wide.pdf}
%   \caption{Wide caption text. \label{fig:wide}}
% \end{figure*}

% Example table:
% \begin{table}[tb]
%   \caption{Table caption. \label{tab:example}}
%   \begin{ruledtabular}
%     \begin{tabular}{lcc}
%       Quantity & Measured & Model \\
%       \colrule
%       $k_r$ & -- & -- \\
%     \end{tabular}
%   \end{ruledtabular}
% \end{table}

% ============================================================
%  IV. CONCLUSION AND PERSPECTIVE
%  - Major conclusion of the project
%  - Proposed future work, and why
% ============================================================
\section{Conclusion and Perspective}
% TODO

% ============================================================
%  ACKNOWLEDGMENTS
% ============================================================
\begin{acknowledgments}
  % TODO: mentor, PI, lab members, Pi2 program / funding.
\end{acknowledgments}

% ============================================================
%  BIBLIOGRAPHY (entries live in references.bib)
% ============================================================
\bibliography{references}

\end{document}
